\documentclass{zkdl-presentation-template}

% Circuit diagrams: inputs gray, gates indigo (+), orange (x), blue (-), wires gray
\colorlet{circle edge}{oc-gray-7}
\colorlet{circle area}{oc-gray-2}
\colorlet{gate1 edge}{oc-indigo-7}
\colorlet{gate1 area}{oc-indigo-1}
\colorlet{gate2 edge}{oc-orange-7}
\colorlet{gate2 area}{oc-orange-1}
\colorlet{gate3 edge}{oc-blue-7}
\colorlet{gate3 area}{oc-blue-1}
\colorlet{wire}{oc-gray-5}
\colorlet{wire out}{oc-gray-7}
\tikzset{
    var/.style={circle, draw=circle edge, fill=circle area, very thick, minimum size=0.7cm, text centered},
    gate1/.style={circle, draw=gate1 edge, fill=gate1 area, ultra thick, minimum size=1cm, text centered},
    gate2/.style={circle, draw=gate2 edge, fill=gate2 area, ultra thick, minimum size=1cm, text centered},
    gate3/.style={circle, draw=gate3 edge, fill=gate3 area, ultra thick, minimum size=1cm, text centered},
    arrow/.style={-Stealth, ultra thick}
}

\title[zk-SNARK: Part I]{\textbf{Introduction to zk-SNARKs. R1CS}}
\author{Blockstream Research}
\date{\today}
\homepage{research.blockstream.com}
\github{BlockstreamResearch/zk-book}

\begin{document}
    \maketitleframe

    \begin{frame}{Plan}
        \tableofcontents
    \end{frame}

    \section{What is zk-SNARK?}

    \begin{frame}{What Is zk-SNARK?}
        \begin{definition}
            \begin{center}
                \textbf{zk-SNARK} \\
                \small\textbf{Z}ero-\textbf{K}nowledge \textbf{S}uccinct \textbf{N}on-interactive \textbf{AR}gument of \textbf{K}nowledge.
            \end{center}
        \end{definition}

        \pause

        \begin{itemize}
            \item \textbf{Argument of Knowledge} --- a proof that the prover knows the data (witness) that resolves a certain
            problem, and this knowledge can be ``extracted''. \pause
            \item \textbf{Succinctness} --- the proof size and verification time is relatively small to the computation size and typically does not depend on the size of 
            the data or statement. \pause
            \item \textbf{Non-interactiveness} --- to produce the proof, the prover does not need any interaction
            with the verifier. \pause
            \item \textbf{Zero-Knowledge} --- the verifier learns nothing about the data used to produce the
            proof, despite knowing that this data resolves the given problem and that the prover possesses it.
        \end{itemize}
    \end{frame}

    \begin{frame}{Still don't get who is SNARK...}
        Well... Let's take a look at some example. \pause

        \begin{columns}
            \begin{column}{0.2\textwidth}
                \includegraphics[width=\linewidth]{images/lecture_8/urtreasurehunter.png}
            \end{column}
    
            \begin{column}{0.7\textwidth}
                Imagine you're part of a treasure hunt...
            \end{column}
        \end{columns}

        \pause

        \begin{columns}
            \begin{column}{0.7\textwidth}
                ...and you've found a hidden treasure chest...
            \end{column}

            \begin{column}{0.2\textwidth}
                \includegraphics[width=\linewidth]{images/lecture_8/uvefoundtreasure.png}
            \end{column}
        \end{columns}

        \pause

        \begin{columns}
            \begin{column}{0.2\textwidth}
                \includegraphics[width=\linewidth]{images/lecture_8/howtomakeaproof.png}
            \end{column}

            \begin{column}{0.7\textwidth}
                ...but how to prove that without revealing the chest location?
            \end{column}
        \end{columns}
    \end{frame}

    \begin{frame}{Still don't get who is SNARK...}
        \begin{columns}
            \begin{column}{0.65\textwidth}
                \begin{block}{The Problem}
                You have found a hidden treasure chest, and you want to 
                prove to the organizer that you know its location without actually revealing that.
                \end{block}
            \end{column}

            \begin{column}{0.3\textwidth}
                \includegraphics[width=\textwidth]{images/lecture_8/treasurelocation.png}
            \end{column}
        \end{columns}

        \pause

        We can retrieve some information from that:

        \vspace{0.1cm}
        \textbf{The Secret Data}: the exact treasure location.

        \vspace{0.1cm}
        \textbf{The Prover}: you.

        \vspace{0.1cm}
        \textbf{The Verifier}: the treasure hunt organizer.
    \end{frame}

    \begin{frame}{Ohh... Got it!}
        Here is how we can apply the zk-SNARK to our problem:

        \begin{itemize} 
            \item \textbf{Argument of Knowledge}: You need to create a proof that demonstrates you know where the
            chest is. \pause
            \item \textbf{Succinct}: The proof you provide is very small and concise. It doesn't matter how
            large the treasure map is or how many steps it took you to find the chest. \pause
            \item \textbf{Non-interactive}: You don't need to have a back-and-forth conversation with the 
            organizer to create this proof. \pause
            \item \textbf{Zero-Knowledge}: The proof doesn't reveal any information about the actual 
            location of the treasure chest.
        \end{itemize}
        \pause
        \vspace{-10pt}
        \begin{columns}
            \begin{column}{0.3\textwidth}
                \includegraphics[width=0.8\textwidth]{images/lecture_8/treasure.png}
            \end{column}

            \begin{column}{0.65\textwidth}
                Well... The golden coin where the pirates' sign is engraved is our zk-SNARK proof!
            \end{column}
        \end{columns}
    \end{frame}

    \begin{frame}{The First Question To Resolve}
        But the problems that we usually want to solve are in a slightly different format.
        \pause

        \vspace{0.1cm}
        When we need to prove that some element is in a merkle tree, we can't come
        to a verifier and give them a "coin"...
        
        \pause
        \begin{questionblock}{Question}
            How do we convert a program into a mathematical language?
        \end{questionblock}

        \center
        \includegraphics[width=8cm]{images/lecture_8/exactcodetomathflow.png}
    \end{frame}

    \section{Boolean Circuits}

    \begin{frame}{Boolean Circuits}
        We can do that in a way like the computer does it --- \textbf{Boolean Circuits}.


        \begin{columns}
            % Description
            \begin{column}{0.7\textwidth}
            \centering
                \begin{minipage}{0.8\textwidth}
                \centering
                % Boolean AND and OR gates
                \begin{tabular}{cc}
                    \begin{tikzpicture}
                        % Nodes
                        \node[var] (a) at (0, -1.5) {$a$};
                        \node[var] (b) at (2, -1.5) {$b$};
                        \node[gate1] (and) at (1, 0) {\texttt{AND}};
                        \node[var] (c) at (1, 1.75) {$c$};
        
                        % Arrows
                        \draw[arrow,wire] (a) -- (and);
                        \draw[arrow,wire] (b) -- (and);
                        \draw[arrow,wire out] (and) -- (c);
                    \end{tikzpicture}
                    &
                    \begin{tikzpicture}
                        % Nodes
                        \node[var] (a) at (0, -1.5) {$a$};
                        \node[var] (b) at (2, -1.5) {$b$};
                        \node[gate2] (or) at (1, 0) {\texttt{OR}};
                        \node[var] (c) at (1, 1.75) {$c$};
        
                        % Arrows
                        \draw[arrow,wire] (a) -- (or);
                        \draw[arrow,wire] (b) -- (or);
                        \draw[arrow,wire out] (or) -- (c);
                    \end{tikzpicture}
                \end{tabular}
                    %\centering
                    %\caption{Boolean \texttt{AND} and \texttt{OR} Gates}
            \end{minipage}
            \end{column}
            \pause
            % Column 2    
            \begin{column}{0.3\textwidth}
                \begin{minipage}{\textwidth}
                \centering
                \begin{tabular}{|c|c|c|}
                    \hline
                    \textbf{A} & \textbf{B} & \textbf{A AND B} \\
                    \hline
                    0 & 0 & 0 \\
                    \hline
                    0 & 1 & 0 \\
                    \hline
                    1 & 0 & 0 \\
                    \hline
                    1 & 1 & 1 \\
                    \hline
                \end{tabular}
                %\caption{\texttt{AND} Gate Truth Table}
            \end{minipage}
        \end{column}
        \end{columns}

        \begin{remarkblock}{Note}
        With any of \{\texttt{AND}, \texttt{NOT}\} or \{\texttt{OR}, \texttt{NOT}\} gates sets one can
        build any possible logical circuit, they are called \textbf{functionally complete} sets.
        \end{remarkblock}
    \end{frame}

    \begin{frame}{Boolean Circuit Example}

        \begin{figure}[h!]
            \centering
            \begin{tikzpicture}
                % Nodes
                \node[var] (a) at (0, -1.5) {$a$};
                \node[var] (b) at (2, -1.5) {$b$};
                \node[var] (c) at (4, -1.5) {$c$};
                \node[gate1] (and) at (1, 0) {\texttt{AND}};
                \node[gate2] (or) at (3, 1.5) {\texttt{OR}};
                \node[var] (d) at (5, 1.5) {$d$};
        
                % Arrows
                \draw[arrow,wire] (a) -- (and);
                \draw[arrow,wire] (b) -- (and);
                \draw[arrow,wire] (c) -- (or);
                \draw[arrow,wire] (and) -- (or);
                \draw[arrow,wire out] (or) -- (d);
            \end{tikzpicture}
            \vspace{-5pt}
            \caption{Example of a circuit evaluating $d = (a \;\texttt{AND}\; b) \,\texttt{OR}\; c$.}
        \end{figure}\pause
        \vspace{-5pt}
        Boolean circuits receive an input vector of $\{0, 1\}$ and resolve to true (1) or false (0);\\
        \vspace{5pt}
        The above circuit can be satisfied with the next values:
        \vspace{-5pt}
        \begin{equation*}
            a = 1, \quad b = 1, \quad c = 0, \quad d = 1
        \end{equation*}
    \end{frame}

    \begin{frame}{SHA-256 Boolean circuit}
        \begin{figure}[h!]
            \centering
            \includegraphics[width=8cm]{images/lecture_8/sha256booleancircuit.jpg}
            \caption{Stats of a SHA256 boolean circuit implementation.}
        \end{figure}
        More than $100000$ gates. Impressive, isn't it? \\
        \vspace{5pt }
        But it also shows how inconvenient the boolean circuits are. 
    \end{frame}

    \section{Arithmetic Circuits}

    \begin{frame}{Arithmetic Circuits}
        Similar to Boolean Circuits, the \textbf{Arithmetic Circuits} consist of gates and
        wires. 
        \begin{itemize}
            \item \textbf{Wires}: elements of some finite field $\mathbb{F}$.
            \item \textbf{Gates}: field addition ($+$) and multiplication ($\times$).
        \end{itemize} 


        \begin{figure}[h!]
            \centering
            % Addition and Multiplication gates
            \begin{tabular}{cc}
                \begin{tikzpicture}
                    % Nodes
                    \node[var] (a) at (0, -1.5) {$a$};
                    \node[var] (b) at (2, -1.5) {$b$};
                    \node[gate1] (add) at (1, 0) {$+$};
                    \node[var] (c) at (1, 1.75) {$c$};

                    % Arrows
                    \draw[arrow,wire] (a) -- (add);
                    \draw[arrow,wire] (b) -- (add);
                    \draw[arrow,wire out] (add) -- (c);
                \end{tikzpicture}
                &
                \begin{tikzpicture}
                    % Nodes
                    \node[var] (a) at (0, -1.5) {$a$};
                    \node[var] (b) at (2, -1.5) {$b$};
                    \node[gate2] (mul) at (1, 0) {$\times$};
                    \node[var] (c) at (1, 1.75) {$c$};

                    % Arrows
                    \draw[arrow,wire] (a) -- (mul);
                    \draw[arrow,wire] (b) -- (mul);
                    \draw[arrow,wire out] (mul) -- (c);
                \end{tikzpicture}
            \end{tabular}
            \caption{Addition and Multiplication Gates}
        \end{figure}
    \end{frame}
 
    \begin{frame}[fragile]{Arithmetic Circuits Example I}
        \begin{example}
            \begin{lstlisting}[language=Python, numbers=none, autogobble=true, xleftmargin=10pt]
                def multiply(a: F, b: F) -> F:
                    return a * b
            \end{lstlisting}
            \pause
            This can be represented as a circuit with only one (multiplication) gate:
            \begin{equation*}
                r = a \times b
            \end{equation*}    
            \pause
            The witness vector (essentially, our solution vector) is $\mathbf{z} = (r, a, b)$, for
            example: $(6, 2, 3)$. \\
            \vspace{5pt}
            We assume that the $a$ and $b$ are input values.
        \end{example}

        \pause

        \begin{remarkblock}{Note}
            We can think of the ``=`` in the gate as an assertion.
        \end{remarkblock}
    \end{frame}

    \begin{frame}[fragile]{Arithmetic Circuits Example II}
        \begin{example}
            Now, suppose we want to implement the evaluation of the polynomial
            $Q(x_1,x_2) = x_1^3 + x_2^2 \in \mathbb{F}[x_1,x_2]$ using arithmetic circuits.
            \begin{lstlisting}[language=Python, numbers=none, autogobble=true, xleftmargin=10pt]
                def evaluate(x1: F, x2: F) -> F:
                    return x1**3 + x2**2
            \end{lstlisting}
            Looks easy, right? \pause But the circuit is now much less trivial.
            \vspace{-8pt}
            \begin{equation*}
                \begin{array}{c c c}
                    \begin{aligned}
                        x_1^2 &= x_1 \times x_1 \\
                        x_1^3 &= x_1^2 \times x_1 \\
                        x_2^2 &= x_2 \times x_2 \\
                        Q &= x_1^3 + x_2^2 
                    \end{aligned}
                    & \text{or} &
                    \begin{aligned}
                        r_1 &= x_1 \times x_1 \\
                        r_2 &= r_1 \times x_1 \\
                        r_3 &= x_2 \times x_2 \\
                        Q &= r_2 + r_3 
                    \end{aligned}
                \end{array}
            \end{equation*}
        \end{example}
    \end{frame}

    \begin{frame}[fragile]{Arithmetic Circuits Example II}
        
        \begin{figure}[h!]
            \centering
            \begin{tikzpicture}
                % Nodes
                \node[var] (x1) at (0, 1) {$x_1$};
                \node[gate2] (x1_x1) at (2, 1) {$\times$};
                \node[gate2] (x1_x1_x1) at (4, 1) {$\times$};
        
                \node[var] (x2) at (0, -1) {$x_2$};
                \node[gate2] (x2_x2) at (2, -1) {$\times$};
        
                \node[gate1] (plus) at (5.0, -0.5) {$+$};
        
                % x1**3
                \draw[arrow,wire] (x1) to [bend left=45] (x1_x1);
                \draw[arrow,wire] (x1) to [bend right=15] (x1_x1);
                \draw[arrow,wire] (x1_x1) -- node[midway, above] {$x_1^2$} (x1_x1_x1);
                \draw[arrow,wire] (x1) to [bend right=45] node[midway, above] {$x_1$} (x1_x1_x1);
        
                % x2**2
                \draw[arrow,wire] (x2) to [bend left=30] (x2_x2);
                \draw[arrow,wire] (x2) to [bend right=30] (x2_x2);
        
                % Summation
                \draw[arrow,wire] (x1_x1_x1) -- node[midway, above right] {$x_1^3$} (plus);
                \draw[arrow,wire] (x2_x2) -- node[midway, above] {$x_2^2$} (plus);
        
                % Result
                \node[var] (q) at (7.0, -0.5) {$Q$};
                \draw[arrow,wire out] (plus) -- (q);
            \end{tikzpicture}
            \caption{Example of a circuit evaluating $x_1^3 + x_2^2$.}
        \end{figure}
    \end{frame}

    \begin{frame}[fragile]{Arithmetic Circuits Example III}
        \begin{example}
            Well, it is quite clear how to represent any polynomial-like expressions. But how can we
            translate \texttt{if} statements?
            \begin{lstlisting}[language=Python, numbers=none, autogobble=true, xleftmargin=8pt]
                def example(a: bool, b: F, c: F) -> F:
                    if a:
                        return b * c 
                    else:
                        return b + c
            \end{lstlisting}
            \pause
            We can transform such a function into the next expression:
            \vspace{-8pt}
            \begin{equation*}
                r = a \times (b \times c) + (1 - a) \times (b + c)    
                \vspace{-8pt}
            \end{equation*}
            \pause
            Corresponding equations for the circuit are:
            \vspace{-8pt}
            \begin{equation*}
                \begin{aligned}
                    r_1 &= b \times c, \quad &r_3 &= 1 - a, \quad &r_5 &= r_3 \times r_2 \\
                    r_2 &= b + c, \quad &r_4 &= a \times r_1, \quad &r &= r_4 + r_5
                \end{aligned}
            \end{equation*}
        \end{example}
    \end{frame}

    \begin{frame}[fragile]{Arithmetic Circuits Example III}

        \begin{figure}[h!]
            \centering
            \begin{tikzpicture}
                % Nodes
                \node[var] (c) at (0.5, -3) {$c$};
                \node[var] (b) at (0.5, -1.5) {$b$};
                \node[var] (a) at (0.5, 0) {$a$};
                \node[var] (one) at (0.5, 1.5) {$1$};
        
                % b+c and b*c gates
                \node[gate1] (b_plus_c) at (3, -1.5) {$+$};
                \node[gate2] (b_times_c) at (3, -3.0) {$\times$};
        
                \node[gate3] (one_minus_a) at (3, 0.75) {$-$};
        
                % a*b*c and (1-a)(b+c) gates
                \node[gate2] (a_times_b_times_c) at (6, -2.0) {$\times$};
                \node[gate2] (one_minus_a_times_b_plus_c) at (6, -0.5) {$\times$};
        
                % a*b*c + (1-a)(b+c) gate
                \node[gate1] (r) at (9, -1.25) {$+$};
        
                % Result node
                \node[var] (result) at (11.5, -1.25) {$r$};
        
                % b+c and b*c arrows
                \draw[arrow,wire] (b) to (b_plus_c);
                \draw[arrow,wire] (b) to (b_times_c);
                \draw[arrow,wire] (c) to (b_plus_c);
                \draw[arrow,wire] (c) to (b_times_c);
        
                % 1 - c arrow
                \draw[arrow,wire] (one) to (one_minus_a);
                \draw[arrow,wire] (a) to (one_minus_a);
        
                % a*b*c and (1-a)(b+c) arrows
                \draw[arrow,wire] (a) to [bend left=20] (a_times_b_times_c);
                \draw[arrow,wire] (b_times_c) to node[midway, above] {$r_1$} (a_times_b_times_c);
                \draw[arrow,wire] (one_minus_a) to node[midway, above] {$r_3$} (one_minus_a_times_b_plus_c);
                \draw[arrow,wire] (b_plus_c) to node[midway, above] {$r_2$} (one_minus_a_times_b_plus_c);
        
                % a*b*c + (1-a)(b+c) arrows
                \draw[arrow,wire] (a_times_b_times_c) to [bend right=20] node[midway, above] {$r_4$} (r);
                \draw[arrow,wire] (one_minus_a_times_b_plus_c) to [bend left=20] node[midway, above] {$r_5$} (r);
        
                % Result arrow
                \draw[arrow,wire out] (r) to (result);
        
            \end{tikzpicture}
            \caption{Example of a circuit evaluating the \texttt{if} statement logic.}
            \label{fig:multivariate-polynomial-circuit}
        \end{figure}
    \end{frame}

    \section{Linear Algebra Recap}

    \begin{frame}{Linear Algebra Recap}
        We will need a few operations from the Linear Algebra lecture (Lecture 3). For $\mathbf{u}, \mathbf{v} \in \mathbb{F}^n$ and matrices $A, B$:
        \begin{itemize}
            \item \textbf{Matrix-vector product:} $A\mathbf{x}$ for $A \in \mathbb{F}^{m \times n}$, $\mathbf{x} \in \mathbb{F}^n$; vectors are columns, $\mathbf{x}^{\top}$ is a row.\pause
            \item \textbf{Dot (inner) product:} $\langle \mathbf{u}, \mathbf{v} \rangle := \mathbf{u}^{\top}\mathbf{v} = \sum_{i=1}^{n} u_iv_i$.\pause
            \item \textbf{Hadamard product:} $(A \odot B)_{i,j} = A_{i,j}B_{i,j}$ (element-wise).\pause
            \item \textbf{Outer product:} $\mathbf{u} \otimes \mathbf{v} := \mathbf{u}\mathbf{v}^{\top}$, a rank-1 matrix for non-zero $\mathbf{u},\mathbf{v}$.
        \end{itemize}
        \pause

        \begin{example}
            For $\mathbf{u} = (1,2,3)$ and $\mathbf{v} = (2,4,3)$: $\langle \mathbf{u}, \mathbf{v} \rangle = 2 + 8 + 9 = 19$ and $\mathbf{u} \odot \mathbf{v} = (2, 8, 9)$.
        \end{example}
    \end{frame}

    \section{Rank-1 Constraint System}

    \begin{frame}{Constraint Definition}
        \begin{definition}
            Each \textbf{constraint} in the Rank-1 Constraint System must be in the form:
            \begin{equation*}
                \langle \boldsymbol{\ell}, \mathbf{z}\rangle \times \langle \boldsymbol{r}, \mathbf{z}\rangle = \langle \boldsymbol{o}, \mathbf{z}\rangle
            \end{equation*}
            Where $\mathbf{z}$ is a vector containing all the \textit{input}, \textit{output}, and 
            \textit{intermediate} variables involved in the computation. The vectors $\boldsymbol{\ell}$, 
            $\boldsymbol{r}$, and $\boldsymbol{o}$ are vectors of coefficients corresponding to these variables, 
            and they define the relationship between the linear combinations of $\mathbf{z}$ on the 
            left-hand side and the right-hand side of the equation.
        \end{definition}        
    \end{frame}

    \begin{frame}{Constraint Example}
        \begin{example}
            Consider the most basic circuit with one multiplication gate: $x_1 \times x_2 = r$.
            The witness vector $\mathbf{z} = (r, x_1, x_2)$. So
            \vspace{-5pt}
            \begin{align*}
                z_2 &\times z_3 = z_1 \\
                (0 + z_2 + 0) &\times (0 + 0 + z_3) = z_1 + 0 + 0 \\
                (0z_1 + 1z_2 + 0z_3) &\times (0z_1 + 0z_2 + 1z_3) = 1z_1 + 0z_2 + 0z_3
            \end{align*}
            Therefore the coefficients vectors are:
            \vspace{-5pt}
            \begin{equation*}
                \boldsymbol{\ell} = (0, 1, 0), \quad \boldsymbol{r} = (0, 0, 1), \quad \boldsymbol{o} = (1, 0, 0). 
                \vspace{-5pt}
            \end{equation*}
            The general form of our constraint is:
            \vspace{-5pt}
            \begin{equation*}
                (\ell_1z_1 + \ell_2z_2 + \ell_3z_3)(r_1z_1 + r_2z_2 + r_3z_3) = o_1z_1 + o_2z_2 + o_3z_3
            \end{equation*}
        \end{example}
    \end{frame}

    \begin{frame}[fragile]{Constraint System Example}
        Now, let us consider a more complex example.
        \begin{lstlisting}[language=Python,numbers=none]
        def r(x1: F, x2: F, x3: F) -> F:
            return x2 * x3 if x1 else x2 + x3
        \end{lstlisting}

        \vspace{-10pt}

        That can be expressed as:
        \begin{equation*}
            r = x_1 \times (x_2 \times x_3) + (1 - x_1) \times (x_2 + x_3)
        \end{equation*}
        
        \pause
        \begin{alertblock}{Note}
            We need a boolean restriction for $x_1$: that is, $x_1 \times (1 - x_1) = 0$.
        \end{alertblock}

        \pause
        Thus, the next constraints can be built:
        \vspace{-5pt}
        \begin{align*}
            x_1 \times x_1 &= x_1 \quad \text{(binary check)} \tag{1} \\
            x_2 \times x_3 &= \mathsf{mult} \tag{2} \\
            x_1 \times \mathsf{mult} &= \mathsf{selectMult} \tag{3} \\
            (1 - x_1) \times (x_2 + x_3) &= r - \mathsf{selectMult} \tag{4}
        \end{align*}
    \end{frame}

    \begin{frame}[fragile]{Constraint System Example}
        The witness vector: $\mathbf{z} = (1, r, x_1, x_2, x_3, \mathsf{mult}, \mathsf{selectMult})$. The coefficients vectors:

        \vspace{-20pt}\small

        \begin{align*}
            \boldsymbol{\ell}_1 &= (0, 0, 1, 0, 0, 0, 0), & \boldsymbol{r}_1 &= (0, 0, 1, 0, 0, 0, 0), & \boldsymbol{o}_1 &= (0, 0, 1, 0, 0, 0, 0) \\
            \boldsymbol{\ell}_2 &= (0, 0, 0, 1, 0, 0, 0), & \boldsymbol{r}_2 &= (0, 0, 0, 0, 1, 0, 0), & \boldsymbol{o}_2 &= (0, 0, 0, 0, 0, 1, 0) \\
            \boldsymbol{\ell}_3 &= (0, 0, 1, 0, 0, 0, 0), & \boldsymbol{r}_3 &= (0, 0, 0, 0, 0, 1, 0), & \boldsymbol{o}_3 &= (0, 0, 0, 0, 0, 0, 1) \\
            \boldsymbol{\ell}_4 &= (1, 0, -1, 0, 0, 0, 0), & \boldsymbol{r}_4 &= (0, 0, 0, 1, 1, 0, 0), & \boldsymbol{o}_4 &= (0, 1, 0, 0, 0, 0, -1)
        \end{align*}\normalsize

        \vspace{-10pt}

        Using the arithmetic in a large $\mathbb{F}_p$, consider the following values:
        
        \vspace{-15pt}
        \begin{equation*}
            x_1 = 1, \quad x_2 = 3, \quad x_3 = 4
            \vspace{-5pt}
        \end{equation*}
    
        Verifying the constraints:
        \begin{enumerate}
            \item \( x_1 \times x_1 = x_1 \quad (1 \times 1 = 1) \)
            \item \( x_2 \times x_3 = \mathsf{mult} \quad (3 \times 4 = 12) \)
            \item \( x_1 \times \mathsf{mult} = \mathsf{selectMult} \quad (1 \times 12 = 12) \)
            \item \( (1 - x_1) \times (x_2 + x_3) = r - \mathsf{selectMult} \quad (0 \times 7 = 12 - 12) \)
        \end{enumerate}
    \end{frame}

    \begin{frame}{R1CS In Matrix Form}
        \begin{theorem} 
            Consider a \textbf{Rank-1 Constraint System} (\textbf{R1CS}) defined over $m$ constraints. Each constraint in such system is
            associated with coefficient vectors $\boldsymbol{\ell}_i$, $\boldsymbol{r}_i$, and $\boldsymbol{o}_i$, where $i \in \{1, 2, \dots, m\}$ and
            also a witness vector $\mathbf{z}$ consisting of $n$ elements.
        
            Then this system can also be represented using the corresponding matrices $L$, $R$, and $O$.
            \begin{align*}
                L = \begin{bmatrix}
                    \ell_{1,1} & \ell_{1,2} & \dots & \ell_{1,n} \\
                    \ell_{2,1} & \ell_{2,2} & \dots & \ell_{2,n} \\
                    \vdots & \vdots & \ddots & \vdots \\
                    \ell_{m,1} & \ell_{m,2} & \dots & \ell_{m,n}
                \end{bmatrix}, \; \text{same for $R$ and $O$,}
            \end{align*}
            such that all constraints can be reduced to the single equation:
            \begin{equation*}
                L\mathbf{z} \odot R\mathbf{z} = O\mathbf{z}
            \end{equation*}
        \end{theorem}
    \end{frame}

    \begin{frame}{R1CS In Matrix Form}
        \textbf{Proof.} Matrices defined this way can be expressed as
        \vspace{-5pt}
        {\scriptsize \begin{equation*}
            L = \begin{bmatrix}
                \boldsymbol{\ell}_1^{\top} \\ \boldsymbol{\ell}_2^{\top} \\ \vdots \\ \boldsymbol{\ell}_m^{\top}
            \end{bmatrix}, \quad R = \begin{bmatrix}
                \boldsymbol{r}_1^{\top} \\ \boldsymbol{r}_2^{\top} \\ \vdots \\ \boldsymbol{r}_m^{\top}
            \end{bmatrix}, \quad O = \begin{bmatrix}
                \boldsymbol{o}_1^{\top} \\ \boldsymbol{o}_2^{\top} \\ \vdots \\ \boldsymbol{o}_m^{\top}
            \end{bmatrix}
            \vspace{-5pt}
        \end{equation*}}
        \pause
        Consider an expression $L\mathbf{z}$:
        \vspace{-5pt}
        {\scriptsize \begin{equation*}
            L\mathbf{z} = \begin{bmatrix}
                \boldsymbol{\ell}_1^{\top} \\ \boldsymbol{\ell}_2^{\top} \\ \vdots \\ \boldsymbol{\ell}_m^{\top}
            \end{bmatrix}\begin{bmatrix}
                z_1 \\ z_2 \\ \vdots \\ z_n
            \end{bmatrix} = \begin{bmatrix}
                \boldsymbol{\ell}_1^{\top}\mathbf{z} \\ \boldsymbol{\ell}_2^{\top}\mathbf{z} \\ \vdots \\ \boldsymbol{\ell}_m^{\top}\mathbf{z}
            \end{bmatrix} = \begin{bmatrix}
                \sum_{i=1}^{n}\ell_{1,i}z_i \\ \sum_{i=1}^{n}\ell_{2,i}z_i \\ \vdots \\ \sum_{i=1}^{n}\ell_{m,i}z_i
            \end{bmatrix} = \begin{bmatrix}
                \langle \boldsymbol{\ell}_1, \mathbf{z}\rangle \\
                \langle \boldsymbol{\ell}_2, \mathbf{z}\rangle \\
                \vdots \\
                \langle \boldsymbol{\ell}_m, \mathbf{z}\rangle 
            \end{bmatrix}
            \vspace{-5pt}
        \end{equation*}}
        \pause
        Therefore, we have:
        \vspace{-5pt}
        {\scriptsize \begin{equation*}
            L\mathbf{z} = \begin{bmatrix}
                \langle \boldsymbol{\ell}_1, \mathbf{z}\rangle \\
                \langle \boldsymbol{\ell}_2, \mathbf{z}\rangle \\
                \vdots \\
                \langle \boldsymbol{\ell}_m, \mathbf{z}\rangle 
            \end{bmatrix}, \quad
            R\mathbf{z} = \begin{bmatrix}
                \langle \boldsymbol{r}_1, \mathbf{z}\rangle \\
                \langle \boldsymbol{r}_2, \mathbf{z}\rangle \\
                \vdots \\
                \langle \boldsymbol{r}_m, \mathbf{z}\rangle 
            \end{bmatrix}, \quad
            O\mathbf{z} = \begin{bmatrix}
                \langle \boldsymbol{o}_1, \mathbf{z}\rangle \\
                \langle \boldsymbol{o}_2, \mathbf{z}\rangle \\
                \vdots \\
                \langle \boldsymbol{o}_m, \mathbf{z}\rangle 
            \end{bmatrix}
            \vspace{-5pt}
        \end{equation*}}
        \pause
        Thus, {\small $L\mathbf{z} \odot R\mathbf{z} = O\mathbf{z}$} is equivalent to the system of $m$ constraints:
        \vspace{-5pt}
        {\small \begin{equation*}
            \langle \boldsymbol{\ell}_j, \mathbf{z}\rangle \times \langle \boldsymbol{r}_j, \mathbf{z} \rangle = \langle \boldsymbol{o}_j, \mathbf{z} \rangle, \; j \in \{1,\dots,m\}.
        \end{equation*}}
    \end{frame}

    \begin{frame}{R1CS In Matrix Form}
        \begin{example}
            The vectors $\boldsymbol{\ell}_i$ from the previous examples have the form:
            {\begin{align*}
                \boldsymbol{\ell}_1 &= (0, 0, 1, 0, 0, 0, 0) \\
                \boldsymbol{\ell}_2 &= (0, 0, 0, 1, 0, 0, 0) \\
                \boldsymbol{\ell}_3 &= (0, 0, 1, 0, 0, 0, 0) \\
                \boldsymbol{\ell}_4 &= (1, 0, -1, 0, 0, 0, 0)
            \end{align*}}
            This corresponds to $n = 7, m = 4$, so the matrix $L$ becomes:
            {\footnotesize \begin{equation*}
                \begin{bmatrix}
                    \ell_{1,1} & \ell_{1,2} & \ell_{1,3} & \ell_{1,4} & \ell_{1,5} & \ell_{1,6} & \ell_{1,7} \\
                    \ell_{2,1} & \ell_{2,2} & \ell_{2,3} & \ell_{2,4} & \ell_{2,5} & \ell_{2,6} & \ell_{2,7} \\
                    \ell_{3,1} & \ell_{3,2} & \ell_{3,3} & \ell_{3,4} & \ell_{3,5} & \ell_{3,6} & \ell_{3,7} \\
                    \ell_{4,1} & \ell_{4,2} & \ell_{4,3} & \ell_{4,4} & \ell_{4,5} & \ell_{4,6} & \ell_{4,7}
                \end{bmatrix} = \begin{bmatrix}
                    0 & 0 & 1 & 0 & 0 & 0 & 0 \\
                    0 & 0 & 0 & 1 & 0 & 0 & 0 \\
                    0 & 0 & 1 & 0 & 0 & 0 & 0 \\
                    1 & 0 & -1 & 0 & 0 & 0 & 0 
                \end{bmatrix}
            \end{equation*}}
        \end{example}
    \end{frame}

    \begin{frame}{Why Rank-1?}
        \begin{lemma}
            Suppose we have a constraint $\langle \boldsymbol{\ell}, \mathbf{z}\rangle \times \langle \boldsymbol{r}, \mathbf{z}\rangle = \langle \boldsymbol{o}, \mathbf{z} \rangle$ 
            with coefficient vectors $\boldsymbol{\ell}$, $\boldsymbol{r}$, $\boldsymbol{o}$ and witness vector 
            $\mathbf{z}$ (all from $\mathbb{F}^n$). Then it can be expressed in the form:
            \vspace{-5pt}
            \begin{equation*}
                 \mathbf{z}^{\top} T \mathbf{z} - \boldsymbol{o}^{\top} \mathbf{z} = 0
                 \vspace{-5pt}
            \end{equation*}
            Where $T = \boldsymbol{\ell} \otimes \boldsymbol{r}$ --- \textbf{rank-1} matrix.
        \end{lemma}
        \vspace{-10pt}
        \pause
        \textbf{Lemma proof.} Consider $\boldsymbol{\ell}, \boldsymbol{r}, \boldsymbol{o}, \mathbf{z} \in \mathbb{F}^n$.
        \vspace{-5pt}
        \begin{equation*}
            \left(\sum_{i=1}^{n} \ell_i z_i\right) \times \left(\sum_{j=1}^{n} r_j z_j\right) = \sum_{k=1}^{n} o_k z_k
            \vspace{-5pt}
        \end{equation*}
        Combine the products into a double sum on the left side:
        \vspace{-5pt}
        \begin{equation*}
            \sum_{i=1}^{n} \sum_{j=1}^{n} \ell_i r_j z_i z_j = \mathbf{z}^{\top} (\boldsymbol{\ell} \otimes \boldsymbol{r}) \mathbf{z} = \mathbf{z}^{\top} T \mathbf{z}
            \vspace{-5pt}
        \end{equation*}
        Thus, the constraint can be written as $\mathbf{z}^{\top} T \mathbf{z} - \boldsymbol{o}^{\top} \mathbf{z} = 0$.
    \end{frame}

    \makeclosingframe
\end{document}